\begin{proof}[Proof of \cref{obj:lem:pinsker-finite}]
\begin{enumerate}
\item Let \(\mathcal Z_{\mathrm{fin}}\) denote the finite alphabet. If some \(z\in\mathcal Z_{\mathrm{fin}}\) satisfies \(\nu'(z)=0\) and \(\nu(z)>0\), then
\[
D(\nu\Vert\nu')=+\infty,
\qquad
\sqrt{D(\nu\Vert\nu')/2}=+\infty,
\]
which proves the inequality in this case. Henceforth assume
\[
\nu'(z)=0\Longrightarrow \nu(z)=0
\qquad(z\in\mathcal Z_{\mathrm{fin}}).
\]
% lean: CausalSmith.Stat.AnnotationRarearmFrontier.PinskerFinite

\item Equip the alphabet with its full power-set sigma-algebra and construct the discrete probability measures
\[
\bar\nu(E)=\sum_{z\in E}\nu(z),
\qquad
\bar\nu'(E)=\sum_{z\in E}\nu'(z)
\qquad(E\subseteq\mathcal Z_{\mathrm{fin}}).
\]
In particular,
\[
\bar\nu(\{z\})=\nu(z),
\qquad
\bar\nu'(\{z\})=\nu'(z).
\]
If \(\bar\nu'(E)=0\), nonnegativity forces \(\nu'(z)=0\) for every \(z\in E\). The support condition then gives
\[
\bar\nu(E)=\sum_{z\in E}\nu(z)=0.
\]
Thus \(\bar\nu\ll\bar\nu'\).

Define the density and log-likelihood functions on the whole alphabet by
\[
\rho_{\mathrm{fin}}(z)=
\begin{cases}
\nu(z)/\nu'(z),&\nu'(z)>0,\\
0,&\nu'(z)=0,
\end{cases}
\qquad
\lambda_{\mathrm{fin}}(z)=
\begin{cases}
\log\rho_{\mathrm{fin}}(z),&\nu(z)>0,\\
0,&\nu(z)=0.
\end{cases}
\]
The singleton density identity is
\[
\rho_{\mathrm{fin}}(z)\bar\nu'(\{z\})
=\bar\nu(\{z\}),
\]
so \(\rho_{\mathrm{fin}}\) is a density of \(\bar\nu\) with respect to \(\bar\nu'\). At every positive \(\nu\)-mass, both masses in the ratio are positive. Consequently, the log-likelihood is integrable on this finite space, and the measure divergence is finite. Writing its real value as
\[
\mathcal K_{\mathrm{fin}}
:=\int\lambda_{\mathrm{fin}}\,d\bar\nu,
\]
evaluation on singletons gives
\[
\mathcal K_{\mathrm{fin}}
=\sum_{z\in\mathcal Z_{\mathrm{fin}}}
  \nu(z)\lambda_{\mathrm{fin}}(z)
=\sum_{z\in\mathcal Z_{\mathrm{fin}}}
  \nu(z)\log\frac{\nu(z)}{\nu'(z)}
=D(\nu\Vert\nu')<\infty.
\]
Here, as in the statement, every zero-\(\nu\)-mass summand is zero.
% lean: Causalean.Mathlib.InformationTheory.klDiv_toReal_eq_sum_measureReal

\item Define the event-supremum distance and the positive-gap event by
\[
\mathsf V_{\mathrm{fin}}
:=\sup_{E\subseteq\mathcal Z_{\mathrm{fin}}}
  |\bar\nu(E)-\bar\nu'(E)|,
\qquad
\mathcal E_{\mathrm{fin},+}
:=\{z\in\mathcal Z_{\mathrm{fin}}:\nu(z)-\nu'(z)\ge0\}.
\]
Normalization yields
\[
\sum_z\bigl(\nu(z)-\nu'(z)\bigr)=0.
\]
Splitting this sum and its absolute-value counterpart over
\(\mathcal E_{\mathrm{fin},+}\) and its complement gives
\[
\begin{aligned}
0
&=\sum_{z\in\mathcal E_{\mathrm{fin},+}}(\nu(z)-\nu'(z))
 +\sum_{z\notin\mathcal E_{\mathrm{fin},+}}(\nu(z)-\nu'(z)),\\
\sum_z|\nu(z)-\nu'(z)|
&=\sum_{z\in\mathcal E_{\mathrm{fin},+}}(\nu(z)-\nu'(z))
 -\sum_{z\notin\mathcal E_{\mathrm{fin},+}}(\nu(z)-\nu'(z))\\
&=2\sum_{z\in\mathcal E_{\mathrm{fin},+}}(\nu(z)-\nu'(z)).
\end{aligned}
\]
The signs defining the event justify the second line. Finite additivity on the disjoint singletons also gives
\[
\bar\nu(\mathcal E_{\mathrm{fin},+})
-\bar\nu'(\mathcal E_{\mathrm{fin},+})
=\sum_{z\in\mathcal E_{\mathrm{fin},+}}(\nu(z)-\nu'(z)).
\]
Every event gap is bounded by the defining supremum; its terms lie in \([0,1]\). Therefore
\[
\operatorname{TV}(\nu,\nu')
=\sum_{z\in\mathcal E_{\mathrm{fin},+}}(\nu(z)-\nu'(z))
\le
\left|\sum_{z\in\mathcal E_{\mathrm{fin},+}}(\nu(z)-\nu'(z))\right|
\le\mathsf V_{\mathrm{fin}}.
\]
% lean: CausalSmith.Stat.AnnotationRarearmFrontier.PinskerFinite.hTV

\item We next bound the absolute density deviation. Define, for \(t_{\mathrm{fin}}\in[0,\infty)\),
\[
\kappa_{\mathrm{fin}}(t_{\mathrm{fin}})
:=t_{\mathrm{fin}}\log t_{\mathrm{fin}}+1-t_{\mathrm{fin}},
\qquad
\Phi_{\mathrm{fin}}(t_{\mathrm{fin}})
:=(t_{\mathrm{fin}}+2)\kappa_{\mathrm{fin}}(t_{\mathrm{fin}})
-\frac32(t_{\mathrm{fin}}-1)^2,
\]
with \(0\log0=0\). The density identity and normalization imply
\[
\int\rho_{\mathrm{fin}}\,d\bar\nu'=1,
\qquad
\mathcal K_{\mathrm{fin}}
=\int\kappa_{\mathrm{fin}}(\rho_{\mathrm{fin}})\,d\bar\nu'.
\]
Indeed, the density identity converts the logarithmic term into
\(\int\lambda_{\mathrm{fin}}\,d\bar\nu\), and the remaining integral is \(1-1=0\).

For \(t_{\mathrm{fin}}>0\), differentiation gives
\[
\begin{aligned}
\Phi_{\mathrm{fin}}'(t_{\mathrm{fin}})
&=2(t_{\mathrm{fin}}+1)\log t_{\mathrm{fin}}
  -4(t_{\mathrm{fin}}-1),\\
\Phi_{\mathrm{fin}}''(t_{\mathrm{fin}})
&=2\bigl(\log t_{\mathrm{fin}}+t_{\mathrm{fin}}^{-1}-1\bigr)
\ge0.
\end{aligned}
\]
The last inequality follows from
\[
-\log t_{\mathrm{fin}}
=\log(t_{\mathrm{fin}}^{-1})
\le t_{\mathrm{fin}}^{-1}-1.
\]
Hence \(\Phi_{\mathrm{fin}}'\) is increasing. Since
\[
\Phi_{\mathrm{fin}}(1)=\Phi_{\mathrm{fin}}'(1)=0,
\]
its derivative is nonpositive on \((0,1]\) and nonnegative on
\([1,\infty)\). Thus \(\Phi_{\mathrm{fin}}\) decreases toward zero on
\((0,1]\) and increases from zero on \([1,\infty)\). At the endpoint,
\(\Phi_{\mathrm{fin}}(0)=1/2\). Dividing its nonnegativity by the positive denominator yields
\[
0\le
\frac{(t_{\mathrm{fin}}-1)^2}{t_{\mathrm{fin}}+2}
\le\frac23\kappa_{\mathrm{fin}}(t_{\mathrm{fin}})
\qquad(t_{\mathrm{fin}}\ge0).
\]

Define the following functions on \(\mathcal Z_{\mathrm{fin}}\):
\[
\begin{aligned}
g_{\mathrm{fin}}(z)
&:=\frac{(\rho_{\mathrm{fin}}(z)-1)^2}{\rho_{\mathrm{fin}}(z)+2},\\
f_{\mathrm{fin},1}(z)
&:=\frac{|\rho_{\mathrm{fin}}(z)-1|}
        {\sqrt{\rho_{\mathrm{fin}}(z)+2}},\\
f_{\mathrm{fin},2}(z)
&:=\sqrt{\rho_{\mathrm{fin}}(z)+2}.
\end{aligned}
\]
All these functions are integrable, together with their squares, on the finite space. The scalar bound gives
\[
0\le\int g_{\mathrm{fin}}\,d\bar\nu'
\le\frac23\mathcal K_{\mathrm{fin}},
\qquad
\mathcal K_{\mathrm{fin}}\ge0.
\]
Moreover,
\[
f_{\mathrm{fin},1}^2=g_{\mathrm{fin}},
\qquad
f_{\mathrm{fin},2}^2=\rho_{\mathrm{fin}}+2,
\qquad
f_{\mathrm{fin},1}f_{\mathrm{fin},2}
=|\rho_{\mathrm{fin}}-1|,
\]
and therefore
\[
\int f_{\mathrm{fin},2}^2\,d\bar\nu'=1+2=3.
\]
Cauchy--Schwarz now yields
\[
\begin{aligned}
\int|\rho_{\mathrm{fin}}-1|\,d\bar\nu'
&\le
\sqrt{\int f_{\mathrm{fin},1}^2\,d\bar\nu'}
\sqrt{\int f_{\mathrm{fin},2}^2\,d\bar\nu'}\\
&=\sqrt{\int g_{\mathrm{fin}}\,d\bar\nu'}\sqrt3\\
&\le\sqrt{\frac23\mathcal K_{\mathrm{fin}}}\sqrt3
=\sqrt{2\mathcal K_{\mathrm{fin}}}.
\end{aligned}
\]
% lean: Causalean.Stat.pinskerBound_of_ac_of_ne_top

\item To relate the event supremum to this integral, define
\[
\delta_{\mathrm{fin}}(z):=\rho_{\mathrm{fin}}(z)-1
\qquad(z\in\mathcal Z_{\mathrm{fin}}).
\]
Then
\[
\int\delta_{\mathrm{fin}}\,d\bar\nu'=0,
\qquad
\bar\nu(E)-\bar\nu'(E)
=\int_E\delta_{\mathrm{fin}}\,d\bar\nu'
\qquad(E\subseteq\mathcal Z_{\mathrm{fin}}).
\]
The first identity follows from normalization; the second follows by integrating the density and subtracting the integral of \(1\). Splitting the zero integral gives
\[
\int_{E^c}\delta_{\mathrm{fin}}\,d\bar\nu'
=-\int_E\delta_{\mathrm{fin}}\,d\bar\nu'.
\]
On each of \(E\) and \(E^c\), integration of
\(\delta_{\mathrm{fin}}\le|\delta_{\mathrm{fin}}|\) and
\(-\delta_{\mathrm{fin}}\le|\delta_{\mathrm{fin}}|\) implies
\[
\left|\int_E\delta_{\mathrm{fin}}\,d\bar\nu'\right|
\le\int_E|\delta_{\mathrm{fin}}|\,d\bar\nu',
\qquad
\left|\int_E\delta_{\mathrm{fin}}\,d\bar\nu'\right|
\le\int_{E^c}|\delta_{\mathrm{fin}}|\,d\bar\nu'.
\]
Adding these inequalities proves
\[
|\bar\nu(E)-\bar\nu'(E)|
\le\frac12\int|\delta_{\mathrm{fin}}|\,d\bar\nu'.
\]
Taking the supremum over events and using the preceding bound, we obtain
\[
\mathsf V_{\mathrm{fin}}
\le\frac12\int|\rho_{\mathrm{fin}}-1|\,d\bar\nu'
\le\frac12\sqrt{2\mathcal K_{\mathrm{fin}}}
=\sqrt{\frac{\mathcal K_{\mathrm{fin}}}{2}}.
\]
The final equality follows by squaring the two nonnegative quantities.
% lean: Causalean.Stat.tvDist_le_half_integral_abs_rnDeriv

\item Combining the finite-alphabet comparison with this bound and the divergence identity gives
\[
\operatorname{TV}(\nu,\nu')
\le\mathsf V_{\mathrm{fin}}
\le\sqrt{\frac{\mathcal K_{\mathrm{fin}}}{2}}
=\sqrt{\frac{D(\nu\Vert\nu')}{2}}.
\]
In the present case, the divergence is finite and nonnegative, so the real inequality carries directly to the extended nonnegative reals, where the square root is the power \(1/2\). Together with the infinite-divergence case, this proves the claim.
% lean: CausalSmith.Stat.AnnotationRarearmFrontier.PinskerFinite
\end{enumerate}
\end{proof}
